\chapter{Stand der Technik}\label{ch:state_of_the_art}
% TODO Reinforcement Learning und Policy Wiederholungen ersetzen, Bilder einsetzen; Änderungen von Waldemar umsetzen
In der heutigen Robotik wird Deep Reinforcement Learning benutzt, um mit klassischer Regelungstechnik schwer modellierbare Verhaltensweisen zu entwickeln. Mehrere Arbeiten (\cite{tangDeepReinforcementLearning2024}, \cite{margolisWalkTheseWays2022}, \cite{heessEmergenceLocomotionBehaviours2017}, \cite{pengLearningLocomotionSkills2017}, \cite{hwangboLearningAgileDynamic2019}) legen nahe, dass Reinforcement Learning in Bereichen der Robotik eine gängige Methode ist, verschiedene Aufgabenstellungen zu bewältigen. Nennenswert ist unter anderem die Bewältigung von Aufgaben, deren Ziele klar definiert sind wie das Greifen und Montage (vgl. \cite{tangDeepReinforcementLearning2024}), welche sowohl in der Simulation als auch in der Praxis auf einem realen Roboter funktioniert (vgl. \cite{hwangboLearningAgileDynamic2019}). % TODO: label bei Glossar

\section{Deep Reinforcement Learning in der Lokomotion}
\subsection{Robuste Lokomotion durch Deep Reinforcement Learning} % TODO Bild von ANYmal
Die Lokomotion ist einer der erfolgreichsten Einsatzbereiche für Deep Reinforcement Learning. Beispielsweise konnten Bewegungsfähigkeiten für den ANYmal, einem vierbeinigen Roboter der ETH Zürich \cite{hutterANYmalHighlyMobile2016}, mithilfe von Deep Reinforcement learning erlernt und anschließend von der Simulation auf die Hardware übertragen wird \cite{hwangboLearningAgileDynamic2019}.

Zur Gelenkregelung werden Aktuatoren eingesetzt, deren Befehle mittels eines neuronalen Netzes berechnet werden \cite{hwangboLearningAgileDynamic2019}. Die Trainingsdatensätze für das Netz bestehen aus realen Roboterdaten, welche sowohl den Verlauf der tatsächlichen Gelenkpositionen im Gegensatz zu angegebenen Positionen als auch die Gelenkgeschwindigkeiten enthält. Das trainierte Netz gibt das Drehmoment aus an den jeweiligen Drehgelenken. Verglichen mit idealisierten oder analytischen Modellen erzielt das Aktuatorennetzes signifikante Simulationsgenauigkeit.

Die Robustheit der erlernten Policies in der Realität wird erhöht, indem die Masse und Schwerpunkte zufällig variiert und die Eingabe der Gelenk- und Körpergeschwindigkeiten mit Rauschen versehen werden \cite{hwangboLearningAgileDynamic2019}. Zudem wurde ein Curriculum implementiert, welches die Schwierigkeit der Aufgaben schrittweise steigert.

Diese Policies wurden ohne hardware-spezifisches oder abgestimmtes Training direkt aus der Simulation auf den realen Roboter übertragen. Diese kann den Befehlen für die Grundgeschwindigkeit genauer und energieeffizienter als der in \cite{hwangboLearningAgileDynamic2019} beschriebene Referenzalgorithmus folgen, was sich in einer geringeren Abweichung von Lineargeschwindigkeit (linear velocity) und Drehgeschwindigkeit der Hochachse (Yaw rate) zeigt. Für hohe Geschwindigkeiten erreicht der ANYmal mit den gelernten Policies eine Geschwindigkeit von $1,5\frac{m}{s}$, was der maximalen Geschwindigkeit des Roboters entspricht. Zusätzlich ist es möglich, dass der Roboter sich aus einer Sturzposition aufrichtet, welche mit früheren Regelungsmethoden kaum möglich war.

\par\vspace{1em}

Während Hwangbo et al. den Schwerpunkt auf die robuste Übertragung einer einzelnen, leistungsfähigen Lokomotion-Policy von der Simulation auf reale Hardware legen, untersuchen Margolis und Agrawal die Frage, ob erlernte Policies nicht nur auf einen Aufgabenbereich anwendbar sind, sondern auch flexibel an unterschiedliche und zuvor unbekannte Umgebungen angepasst werden kann.

\subsection{Dynamische und vielseitige Fortbewegungsfähigkeiten für Lokomotion}

\subsubsection{End-to-End Ansatz für die Fortbewegungsfähigkeit}
In \cite{margolisWalkTheseWays2022} zeigen Margolis und Agrawal, dass eine verbesserte Generalisierung von gelernten Policies möglich und damit die Effektivität in bislang unbekannten Umgebungen erhöht werden kann. Das von den Autoren als Multiplicity of Behavior bezeichnete Konzept beschreibt eine Policy, welche Trainingsaufgaben auf unterschiedliche Weise lösen kann. Dies ist durch ein Training in einer vereinfachten Umgebung möglich, welche auf unbekannte, dynamische Umgebungen übertragbar ist, um ein Neutraining für spezifische Probleme zu vermeiden.

Die Anpassung der Policy an unbekannte Umgebungen erfolgt durch einen Menschen im laufenden Betrieb, der einen achtdimensionalen Vektor $b_t$ für Verhaltensparameter steuert, um eine passende Laufstrategie zu finden \cite{margolisWalkTheseWays2022}. Mit dem Vektor werden Gangmuster, Körperhaltung und Fußbewegungen an das jeweilige Problem bzw. Terrain angepasst, und können so unter anderem das Überwinden von Treppen oder das Kriechen unter Hindernissen lernen.

\par\vspace{1em}

Diese Lernfähigkeit stellt einen wichtigen Schritt in Richtung robuster, generaliserungsfähiger Lokomotion dar. Gleichzeitig argumentiert Recht, dass eine direkte Steuerung aller Freiheitsgrade durch Deep Reinforcement Learning oft mit Herausforderungen hinsichtlich Stabilität und Sicherheit verbunden ist \cite{rechtTourReinforcementLearning2018}. Aus diesem Grund untersuchen zahlreiche Arbeiten hybride Ansätze, bei denen Deep Reinforcement Learning für hochrangige Entscheidungs- und Verhaltensparameter eingesetzt wird, während die eigentliche Körper- und Aktuatorregelung klassisch geregelt wird.

\subsubsection{Hybrider Regelungsansatz für die Fortbewegungsfähigkeit}
Yin et al. nutzen in \cite{yinRunDogLearning2021} diesen hybriden Ansatz, bei dem die Ansteuerung über eine sog. Whole-Body Control (WBC), sprich der Ganzkörperregelung, erfolgt. Durch die Berechnung der Drehmomente wird eine präzise Bewegungsverfolgung garantiert. Um zusätzlich noch das Gangmuster natürlicher zu gestalten, werden mit Deep Reinforcement Learning die Parameter für den Bewegungsstil generiert und der WBC zugeführt. 

Das Training versucht aus gegebenen Motion-Capture-Clips (MoCap-Clips) von echten, vierbeinigen Tieren die Parameter für das gegebene Gangmuster zu extrahieren \cite{yinRunDogLearning2021}, und dabei anhand dieser Daten den Gang zu imitieren. Die Bewertung erfolgt hier durch einen Vergleich der Geschwindigkeit von den Füßen zwischen den erlernten Bewegungen und den MoCap-Clips, wobei eine hohe zeitliche Übereinstimmung erreicht wird.

\par\vspace{1em}

Während diese Fortbewegungsfähigkeit auf dem Boden von vielen Arbeiten untersucht werden, existieren Problemstellungen, bei denen diese Anpassungsversuche auf ihre Grenzen stößt. Vor allem ein Höhenunterschied oder ein Hindernis, welches nicht durch Klettern überwindbar ist, kann durch einen gezielten Sprung überwunden werden. Bussola et al. untersuchen diese Aufgabenstellung und bedienen sich dabei ebenfalls eines hybriden Ansatzes \cite{bussolaGuidedReinforcementLearning2025}.
 
\subsubsection{Sprungfähigkeit für einen vierbeinigen Roboter}
Die Arbeit von Bussola et al. stellt die Geh und Laufmanöver im Gegensatz zum Sprungmanöver, da die beim Laufen und Gehen der Roboter kontinuierlich mit den Boden interagiert, während beim Sprung die Sprungtrajektorie nach dem Absprung nicht mehr geändert werden kann \cite{bussolaGuidedReinforcementLearning2025}.
Zusätzlich kommt die Landungsschwierigkeit hinzu, da eine Landung den Aufprall abfangen und das Umkippen des Roboters verhindert werden soll.

Neben den Schwierigkeiten, die mit der beim Springen und Landen mit einhergehen, hat der Roboter inhärente Steuerungsbeschränkungen. Die Optimierung des Absprungs muss unter anderem die dynamische Beschränkung der Gelenkposition, Gelenkgeschwindigkeits- und Drehmomentsgrenzen berücksichtigen, die Empfindlichkeit der Trajektorienplanung gegenüber Fehlern und die beschränkte Echtzeitberechnung, die bei einer hohen Frequenz kontinuierlich durchgeführt werden muss und umfangreiches Wissen über das System und dessen Parameter erfordert \cite{bussolaGuidedReinforcementLearning2025}. Zusätzlich stellt die Bodenbeschaffenheit beim Springen und Landen eine Schwierigkeit dar, da dies abhängig vom Boden unterschiedliche Anforderungen an den Sprung bzw. Landung 

Ähnlich wie von Recht \cite{rechtTourReinforcementLearning2018} und Henderson et al. \cite{hendersonDeepReinforcementLearning2019} festgestellt, weisen Bussola et al. darauf hin, dass reine End-to-End Deep Reinforcement Learning Ansätze zwar eine mögliche Alternative darstellen, allerdings hohe Varianz aufweisen, welches die Sicherheit negativ beeinträchtigt.\newline
Um dieses Problem zu umgehen, stellen Bussola et al. den sog. Guided, also geführten, Reinforcement Learning Ansatz vor, bei dem die gesamte Sprungaktion – also von Stillstand bis Landung – als einen einzigen Aktionsschritt geplant wird. Die Sprungtrajektorie wird dabei mithilfe von Bézier-Kurven und einem Uniformly Accelerated Rectilinear Motion (UARM) Modell parametrisiert \cite{bussolaGuidedReinforcementLearning2025}. \newline
Bézier-Kurven beschreiben den Verlauf der Trajektorie über eine kleine Anzahl an Kontrollpunkten, wohingegen das UARM-Modell eine vereinfachte physikalische Beschreibung der Schwerpunktbewegung während des Sprungs beschreibt.

Diese Parametrisierung wird mithilfe von PPO durchgeführt, und versucht dabei, Parameter zu optimieren, welche den Werten aus der optimalen Steuerung ähneln \cite{bussolaGuidedReinforcementLearning2025}.
Die Belohnungsfunktion im Training stellt dabei sicher, dass die Policy die physikalischen und Aktuatorengrenzen während des Absprungs respektiert.

Im Gegensatz zu reinen End-to-End Ansätzen ist die Ausführung der Bewegung vom Lernprozess entkoppelt. Die optimierten Trajektorienparameter werden hierbei an einen klassischen Low-Level-Regler übergeben, der die Bewegung stabil und mit Rücksicht auf die Systemdynamik ausführt.

\subsection{Curriculum Learning für komplexe Aufgabenbewältigung}
Curriculum Learning ist eine Trainingsstrategie im Reinforcement Learning, bei der ein Agent eine explizit entworfene Sequenz von Aufgaben mit steigendem Schwierigkeitsgrad durchläuft, um die Lerngeschwindigkeit und die Endleistung bei einer komplexen Zielaufgabe zu verbessern \cite{francois-lavetIntroductionDeepReinforcement2018, arulkumaranDRL2017}. Im Gegensatz zu einer gleichzeitigen starken zufälligen Verteilung aller Umgebungsparameter, wie sie bei Domain Randomization eingesetzt wird, erlaubt ein Curriculum eine kontrollierte Progression der Schwierigkeit. Dadurch wird die Exploration strukturiert geführt und instabile Trainingsphasen in hochkomplexen Zustandsräumen werden vermieden. In der Lokomotion von Robotern wird dieser Ansatz insbesondere genutzt, um anspruchsvolle Fähigkeiten wie das Überwinden von Hindernissen, Treppensteigen oder Parkour-Manöver effizient zu erlernen, indem der Roboter zunächst auf einfachem Gelände trainiert wird, bevor Parameter wie Stufenhöhe, Steigungswinkel oder Grabenbreite schrittweise erhöht werden \cite{heessEmergenceLocomotionBehaviours2017, rudin2022learningwalkminutesusing, hoellerANYmalParkourLearning2023, jinResilientLeggedLocal2023}.

So zeigen Heess et al. \cite{heessEmergenceLocomotionBehaviours2017}, dass komplexe Verhaltensweisen wie Springen oder Ducken aus einfachen Belohnungssignalen wie Vorwärtsgeschwindigkeit resultieren können, wenn die Umgebungsvariabilität schrittweise erhöht wird. Ein explizites Terrain-Curriculum, also Gelände mit steigendem Schwierigkeitsgrad, verbessert hierbei die Lerngeschwindkeit deutlich \cite{heessEmergenceLocomotionBehaviours2017}.

In der Arbeit von Rudin et al. \cite{rudin2022learningwalkminutesusing} wird, ähnlich wie Heess et al. \cite{heessEmergenceLocomotionBehaviours2017} zeigen, ebenfalls auf verschiedenen Terrain-Levels trainiert, bei dem ein Roboter in schwierigere Level befördert wird bei Bewältigung des Terrains. Nach der Überquerung des schwersten Levels wird der Roboter zufällig zurück versetzt, um katastrophales Vergessen der gelernten Policy zu vermeide \cite{rudin2022learningwalkminutesusing}. Für diesen Prozess setzen Rudin et al. auf eine GPU-beschleunigte, massiv parallele Simulation, die tausende Instanzen der Trainingsumgebungen gleichzeitig ausführt und dabei kein manuelles Tuning der Schwierigkeitsstufen erfordert \cite{rudin2022learningwalkminutesusing}. \newline % TODO Bild von Curriculum

Während die bisher betrachteten Arbeiten den Schwerpunkt auf die Generierung robuster Fortbewegungsmuster und deren Generalisierung konzentrieren, rückt in komplexen Anwendungsszenarien zunehmend die Frage nach der zielgerichteten Bewegung im Raum in den Vordergrund. Neben der Stabilität einzelner Bewegungen ist entscheidend, wie ein Roboter unter Berücksichtigung von Umweltinformationen, Hindernissen und Zielvorgaben geeignete Bewegungsentscheidungen trifft.

Eine enge Verzahnung von Wahrnehmungen, Entscheidungsfindung und Bewegungsausführung zeigen Jin et al. \cite{jinResilientLeggedLocal2023}. Die Autoren trainieren einen End-to-End Navigationsansatz, bei dem rohe, teilweise verrauschte Sensordaten direkt in Navigationskommandos überführt und damit robuste Fortbewegung mit hochrangiger Entscheidungsfindung verbindet.

Diese Aufgabenstellung wird in der Robotik unter dem Begriff Navigation zusammengefasst.

\section{Navigation mithilfe von Deep Reinforcement Learning}
In ihrer Arbeit stellen Tang et al. fest, dass Deep Reinforcement Learning in Bereichen wie das autonome Fahren in Städten oder den Drohnenflug im Freien wenig Verbreitung gefunden haben, unter anderem aufgrund strenger Robustheitsanforderungen sowie fehlender Generalisierbarkeit und Erklärbarkeit \cite{tangDeepReinforcementLearning2024}. Hinzu kommt, dass Navigation Planung mit langer Voraussicht erfordert und währenddessen komplexe Bewegungsdynamiken berücksichtigt werden müssen.

Zusätzlich dazu steht die Lokomotion vierbeiniger Roboter vor der Herausforderung, schwieriges und unstrukturiertes Terrain zu überwinden. Tang et al. berichten in \cite{tangDeepReinforcementLearning2024} unter Verweis auf aktuelle Forschungen , dass zur Erleichterung des Sim-to-Real Transfer ebenfalls ein hybrider Ansatz aus klassischer Regelung und Navigationstraining verwendet wird. Dadurch kann sich das Lernen in der Simulation auf die Verarbeitung visueller Sensordaten und die hochrangige Navigationsentscheidung konzentrieren.

Jin et al. \cite{jinResilientLeggedLocal2023} selbst verfolgen einen End-to-End Ansatz, bei dem aus rohen, teilweise verrauschten Sensordaten direkt Navigationskommandos erzeugt werden, während eine vortrainierte Low-Level-Policy die stabile Lokomotion sicherstellt. Das Training erfolgt mittels eines Curriculums, bei dem Rauheit, Stufenhöhe sowie Anzahl von Hindernissen und Gruben adaptiv erhöht werden. Zusätzlich besitzt der Roboter Information über die eigene Basis-Geschwindigkeit, Verlauf der Beschleunigung, Schwerkraftsvektor, Gelenkposition-/geschwindigkeiten und -drehmomente, mit diesen der Roboter Kollisionen erfassen kann und Bewegungsentscheidungen angepasst werden können \cite{jinResilientLeggedLocal2023}. % TODO Bild von jinResilientLeggedLocal2023

% TODO Generalisierung beim Lernen durch "Ground Truth" (zufügen ins Glossar)
Diese Informationen sind nicht gleich auf die beiden Netzwerke, Actor und Critic, verteilt. Der Actor erhält nur die korrupten und unvollständigen Daten, während der Critic Zugriff auf alle Informationen der Umgebung und des Roboters, dem "Ground Truth", erhält \cite{jinResilientLeggedLocal2023}. Während des Trainings wird dann versucht, die Merkmale der Umgebungen so nahe wie möglich am "Ground Truth" vom Critic abzubilden, um gelernte Policies zu generalisieren und so nicht mehr von der fehlerhaften Karte abhängig zu sein \cite{jinResilientLeggedLocal2023}. Das führt zum Erlangen von Navigationsfähigkeiten, bei der für den Roboter unsichtbare Hindernisse nach Kollision umsteuert werden und es sogar möglich ist, sich aus unsichtbaren Gruben raus zu ziehen \cite{jinResilientLeggedLocal2023}. 

Bei der Auswertung stellen Jin et al. fest, dass die Policy sogar bei 0\% Sichtbarkeit die Hindernisse mit einer Erfolgsrate von über 80\% umgeht, und damit klassische Planungsalgorithmen sowohl beim Bewältigen des Terrains als auch bei der Reduzierung von Kollisionen übertrifft; die Evaluation erfolgt dabei nicht nur in Simulation, sondern auch auf dem realen ANYmal \cite{jinResilientLeggedLocal2023}.

\section{Reward Shaping in Deep Reinforcement Learning} \label{sec:reward_shaping} 
Reward Shaping wird in der Robotik als zentrale Methode eingesetzt, um die Ineffizienz von sogenannten Sparse Rewards, also eher spärlich verteilten Belohnungen, in hochdimensionalen Aktionsräumen zu überwinden \cite{burda2018explorationrandomnetworkdistillation}. \newline 
Bei der Lokomotion von Laufrobotern haben sich sog. Barrier-Based Style Rewards als Ansatz etabliert, bei denen logarithmische Barrierefunktionen genutzt werden, um gewünschte Gangarten oder Körperhöhen als weiche Randbedingungen zu integrieren \cite{kimLearningFrameworkDiverse2025}. Diese Form von Reward Shaping ermöglicht eine präzise Führung des Roboters zu einem natürlichen Bewegungsbild, im Gegensatz zu Constrained RL, welche eine harte Einhaltung der Randbedingung durchsetzt, und damit üblicherweise Instabilität bzw. unendliche Werte bei Randverletzungen hervorruft \cite{kimLearningFrameworkDiverse2025}.

Daneben wird das Risiko des Reward Hackings adressiert, welches Schwachstellen der Belohnungsfunktion ausnutzt, um hohe Belohnungserträge zu erzielen, ohne die eigentlich beabsichtigte Aufgabe korrekt zu lösen \cite{skalseDefiningCharacterizingReward2025}. Frameworks wie ORSO für die automatisierte Belohnungsauswahl \cite{zhangORSOAcceleratingReward2025} oder das EST-Framework für Invarianz-basierte Stress-Tests \cite{shihabDetectingProxyGaming2026} gewinnen hierbei immer an Bedeutung, und stellen sicher, dass Leistungssteigerung auf tatsächlichem Fortschritt beruhen und nicht auf der Ausnutzung der Schwachstellen der Bewertungsfunktionen \cite{zhangORSOAcceleratingReward2025}. \newline
ORSO beschleunigt den Prozess des Belohnungsdesign proaktiv, indem es mittels LLMs Belohnungskandidaten generiert und durch adaptive Bandit-Algorithmen die effizientesten Reward Funktionen auswählt \cite{zhangORSOAcceleratingReward2025}. Im Gegensatz dazu nutzt das EST-Framework invarianzbasierte Stresstests, welche die Sensitivität der Belohnung gegenüber aufgabenirrelevanten Umgebungsdetails messen, um Reward Hacking diagnostisch zu erkennen \cite{shihabDetectingProxyGaming2026}. % TODO LLM im Glossar erklären

\subsection{Dense und Sparse Rewards}
Obwohl Reward Shaping und das damit verbundene Risiko des Reward Hacking die methodische Gestaltung beeinflussen, bleibt die grundlegene Herausforderung der Belohnungsdichte bestehen. Während Reinforcement learning bei dense rewards, wie sie etwa in der Lokomotion durch die Abweichung von einer Soll-Geschwindigkeit natürlich vorkommen und effizient konvergiert, scheitern klassische Algorithmen oft an sparse rewards  \cite{burda2018explorationrandomnetworkdistillation, tangDeepReinforcementLearning2024}. \newline
In diesen Szenarien erhält der Agent nur selten ein nicht-triviales Signal, was zum sog. "Plateau-Problem" führt \cite{burda2018explorationrandomnetworkdistillation}, das bereits grundlegend in \cite{sutton2018reinforcement} diskutiert wird und die ziellose Exploration im Zustandsraum ohne unmittelbare Hinweise für lange Zeit beschreibt. 

Moderne Frameworks wie HERO (\cite{taoHybridReinforcementWhen2025}) verfolgen daher einen hybriden Ansatz, bei der spärlich verteilte, binäre Rewards mit dicht verteilten, kontinuierlichen Scores von Belohnungsmodellen kombiniert wird. Durch stratifizerite Normalisierung wird sichergestellt, dass dichte Signale den Lernprozess innerhalb der durch die spärlich verteilten Signale definierten Grenzen verfeinern, ohne die Korrektheit des Gesamtziels zu gefährden \cite{taoHybridReinforcementWhen2025}. Zusätzlich wird durch das hybride Design ein kontinuierlicher Gradientenfluss der Belohnungen selbst in Sättigungsbereichen ermöglicht, in denen rein binäre Belastungsfunktionen keine verwertbare Lernsignale mehr liefern würden \cite{taoHybridReinforcementWhen2025}.

\subsection{Overfitting in Deep Reinforcement Learning}
Die Optimierung komplexer Belastungsfunktionen führt unmittelbar zur Frage der Generalisierungsfähigkeit, da Deep Reinforcement Learning Systeme aufgrund ihrer hohen Modellkapazität und langen Trainingszeiten stark zu Overfitting neigt \cite{zhang2018studyoverfittingdeepreinforcement}. Dabei wird anstatt eine robuste Strategie zu lernen, spezifische Merkmale und Aktionssequenzen der Trainingsumgebung von den Agenten lediglich auswendig gelernt \cite{zhang2018studyoverfittingdeepreinforcement, sutton2018reinforcement}. Oft tritt die gelernte Policy \say{robust} auf, was bedeutet, dass Overfitting nicht verhindert oder sichtbar gemacht wird trotz Einführung von Stochastizität, z.B. durch zufällige Startpositionen \cite{zhang2018studyoverfittingdeepreinforcement}.

Häufig manifestiert sich Overfitting in der Robotik als Simulation Bias, bei dem die physikalische Ungenauigkeiten der Simulationsumgebung genutzt wird, weshalb eine Portierung auf den realen Roboter scheitert \cite{Kober_IJRR_2013}. Zur Vermeidung werden neben Regularisierungstechniken wie L2-Regularisierung, Dropout oder Batch-Normalisierung zunehmend explizit getrennte, oft przedural generierte Trainings- und Testumgebungen eingesetzt \cite{cobbe2019quantifyinggeneralizationreinforcementlearning, kolomeytsevHybridMotionPlanning2025}.


